\documentclass[11pt]{article}

\usepackage[margin=1in]{geometry}
\usepackage{booktabs}
\usepackage{siunitx}

\title{Assignment 1: Branch Predictor Implementation}
\author{Course 6385}
\date{\today}

\begin{document}
\maketitle

\section{What was implemented}

The predictor is the class \texttt{pm\_predictor} in
\texttt{cbp2-infrastructure-v2/src/my\_predictor.h}. It is the simplified
hybrid from Appendix~B of the assignment: a bimodal table and a two-way
global table. The assignment allows a two-way global predictor in place of
the Pentium~M's four-way table.

A conditional branch is predicted as follows.

\begin{itemize}
  \item The bimodal table has 16 two-bit saturating counters. The index is
        \texttt{PC[5:2]}. The prediction is the counter's high bit.
  \item A 4-bit pattern history register records the outcomes of conditional
        branches. This is the global history register in Appendix~B. It is
        shifted after the branch resolves, and the new taken or not-taken bit
        is inserted.
  \item Those four history bits are XORed with \texttt{PC[5:2]}. Bits
        \texttt{[3:2]} of the result select one of four sets. Bits
        \texttt{[1:0]} are the tag.
  \item Each set has two ways. A way hits when it is valid and the tag
        matches. On a hit the global counter's high bit is the prediction.
        On a miss the bimodal prediction is used.
  \item The global entry is updated either way. A miss fills an invalid way,
        or replaces the least-recently-used way if both are valid. The chosen
        way becomes most recently used.
\end{itemize}

Counters saturate at 0 and 3. The initial state matches the Appendix~B tables:
bimodal counter \(i\) starts at \(i \bmod 4\), and the global ways start with
the valid bits clear and the counter and LRU values shown in that example.

The pattern history register is the simple stand-in for the Pentium~M path
information register. Both are XORed with the branch address to index the
global table. The path register stores address and target bits and is updated
only on taken and indirect branches. This design stores one outcome bit per
conditional branch.

\section{What was left out}

Loop prediction, the branch target buffer, and the path information register
are not implemented. Unconditional branches are predicted taken and do not
update the tables. Target addresses are not predicted; the simulator still
reports direction mispredictions per kilo-instruction, which is the metric
for this assignment.

\section{Check against Appendix~B}

Appendix~B gives 22 conditional branches (\texttt{b1} at \texttt{0x44},
\texttt{b2} at \texttt{0x6C}, \texttt{b3} at \texttt{0x90}) and the expected
state of both tables. Those tables are stored as \texttt{hybrid.csv}, \texttt{local.csv}, and
\texttt{global.csv}. The program
\texttt{appendix\_b\_check} feeds each program counter and outcome to
\texttt{pm\_predictor} and compares the history register, the XOR index, the
16 bimodal counters, and every global valid bit, tag, counter, and LRU bit.

The run on the current predictor reports 22 of 22 rows matching.

\section{Trace results}

The predictor was built with \texttt{make} in
\texttt{cbp2-infrastructure-v2/src} and run on the 20 CBP2 traces. Each trace
is 100 million instructions. MPKI is \(1000 \times (\textit{direction
mispredictions}) / 10^{8}\), counted only for conditional branches. The
always-taken column is the same traces run with the original stub, which
predicted every conditional branch taken.

\begin{table}[ht]
\centering
\begin{tabular}{l S[table-format=3.3] S[table-format=2.3]}
\toprule
{Trace} & {Always taken} & {Hybrid} \\
\midrule
164.gzip      & 70.025 & 49.708 \\
175.vpr       & 85.312 & 40.103 \\
176.gcc       & 86.071 & 61.200 \\
181.mcf       & 102.497 & 41.210 \\
186.crafty    & 48.651 & 40.400 \\
197.parser    & 78.421 & 38.632 \\
201.compress  & 76.983 & 26.885 \\
202.jess      & 92.343 & 47.089 \\
205.raytrace  & 109.506 & 11.827 \\
209.db        & 91.451 & 32.173 \\
213.javac     & 109.644 & 13.131 \\
222.mpegaudio & 110.339 & 12.785 \\
227.mtrt      & 112.022 & 12.275 \\
228.jack      & 77.958 & 41.810 \\
252.eon       & 24.226 & 22.745 \\
253.perlbmk   & 66.833 & 57.871 \\
254.gap       & 75.950 & 44.542 \\
255.vortex    & 70.524 & 38.671 \\
256.bzip2     & 178.257 & 77.119 \\
300.twolf     & 59.221 & 45.328 \\
\midrule
{Average}     & 86.311 & 37.775 \\
\bottomrule
\end{tabular}
\caption{Direction mispredictions per kilo-instruction.}
\end{table}

The hybrid predictor is lower than always-taken on every trace. The average
falls from 86.311 to 37.775 MPKI. The largest drops are on raytrace, javac,
mpegaudio, and mtrt, which are heavily biased and fit even a small bimodal
table. gcc, perlbmk, and bzip2 stay high. The global table has only four sets
and two tags, so unrelated branches alias, and a hit can override a correct
bimodal prediction with a counter trained by a different branch.

\section{Conclusion}

The submitted predictor is the Appendix~B hybrid: a 16-entry bimodal table, a
4-bit pattern history register, and a 2-way tagged global table. A global hit
overrides the bimodal prediction. The implementation matches the worked
example cell for cell, and on the CBP2 traces it cuts the always-taken miss
rate from 86.311 to 37.775~MPKI.

\end{document}
